\documentclass[11pt]{article}
\usepackage[margin=1in]{geometry}
\usepackage{booktabs}
\usepackage{array}
\usepackage{enumitem}
\usepackage[colorlinks=true,urlcolor=blue,linkcolor=black]{hyperref}
\setlist{itemsep=2pt,topsep=3pt}

\title{Wearables and Habits: Project Plan Proposal\\\large Issue \#1 --- for discussion}
\date{3 October 2026}
\author{}

\begin{document}
\maketitle

\noindent\textbf{Status:} proposal, nothing here has been tested. Statements about the dataset are from memory of the
published description and must be confirmed in the data audit.

\section{Summary of the proposal}
\textbf{Question:} do wearables (fitness trackers, GPS watches) causally increase running?
\textbf{Data:} Open e-commerce 1.0 (Berke et al., 2024; DOI \href{https://doi.org/10.7910/DVN/YGLYDY}{10.7910/DVN/YGLYDY}),
which, as I understand it, contains Amazon purchase histories linked to survey demographics.
\textbf{Design:} staggered-adoption event study, with the first wearable purchase as treatment.

\textbf{Central caveat.} The data observe \emph{purchases}, not activity. Running can only be measured through proxies
(e.g.\ running shoes, apparel, gear). The research question should therefore be stated as the effect of wearable
adoption on \emph{running-related purchasing}, with the link to actual running treated as an assumption, not a result.
The second threat is selection: people who decide to start running buy a watch and shoes together, which breaks
no-anticipation. The plan is built around these two issues.

\section{Questions the plan must answer}

\subsection*{Scope}
\begin{itemize}
  \item \textbf{Treatment:} first purchase of a wearable (tracker, GPS/sports watch). Smartwatches are a borderline
        class; proposal is to start with a narrow definition and treat the broad one as robustness.
  \item \textbf{Primary outcome:} running-related purchases (shoes, apparel, accessories) per user-period.
        \textbf{Secondary:} replacement cadence of running shoes as a volume proxy; other fitness categories.
  \item \textbf{Population:} users with an observed wearable purchase (treated) and users with none (never-treated),
        restricted to users with enough purchase history before and after adoption.
  \item \textbf{Time window:} the full span of the dataset, trimmed by the audit. Out of scope: health outcomes,
        non-purchase behavior, anything requiring activity data.
\end{itemize}

\subsection*{Data: sufficient, and what else?}
\begin{itemize}
  \item \textbf{Strengths:} individual purchase histories with timestamps and demographics permit event-study
        timing and heterogeneity by age, sex, income.
  \item \textbf{Weaknesses:} purchases are not use; gifts and household purchases; a single retailer;
        possibly small treated sample once cohorts are split.
  \item \textbf{Supplements to scope, not to commit to:} aggregate public activity data (e.g.\ platform
        trend reports) as external validity for the proxy; a short validation using any source that links
        purchases to self-reported running (the dataset's own survey, if it has such items).
  \item \textbf{Alternatives if the audit fails:} different treatment/outcome pair in the same data
        (e.g.\ other fitness equipment), or a different dataset with activity logs but weaker identification.
\end{itemize}

\subsection*{Identification}
\begin{tabular}{p{0.28\textwidth}p{0.66\textwidth}}
\toprule
Element & Proposal \\
\midrule
Estimator & Heterogeneity-robust staggered DiD (Callaway--Sant'Anna as primary; Sun--Abraham or
            de~Chaisemartin--D'Haultfoeuille as cross-check). Avoid plain two-way fixed effects. \\
Comparison group & Never-treated users; not-yet-treated as robustness. \\
Key assumptions & Parallel trends; no anticipation; stable composition of users. \\
Tests & Pre-trend event-study plots; placebo outcomes (unrelated categories); placebo treatment
        (other electronics purchases); sensitivity of trends (HonestDiD-style bounds);
        dropping users whose running-gear purchases coincide with the wearable within a short window. \\
Alternatives & Matched DiD on pre-period purchases; heterogeneity cuts. IV is not proposed (no credible instrument known). \\
\bottomrule
\end{tabular}

\subsection*{Feasibility checks before committing}
\begin{enumerate}
  \item Number of treated users, and how many treatment cohorts (by period).
  \item Pre- and post-adoption history length per treated user.
  \item Share of users with any running-related purchase; how sparse the outcome is.
  \item Whether wearables and running categories can be classified reliably from the product fields.
  \item Minimum detectable effect from a simulation using the real treatment timing and outcome variance.
\end{enumerate}

\subsection*{Outputs}
A LaTeX report per phase (compiled PDF) under \texttt{reports/}; scripts under \texttt{code/scripts/};
figures and tables under \texttt{code/analysis/outputs/}; a final summary report. A paper in \texttt{paper/} is
optional and decided at the go/no-go point.

\section{Phase outline}
\begin{tabular}{p{0.04\textwidth}p{0.18\textwidth}p{0.68\textwidth}}
\toprule
\# & Phase & Content (headings only) \\
\midrule
1 & Setup & Container, repo structure, data download with retry/resume, data dictionary. \\
2 & Data audit & Classification of wearables and running products, feasibility checks, power. \\
\multicolumn{3}{l}{\textbf{GO/NO-GO:} after phase 2 (see below).} \\
3 & Design & Fix the treatment and outcome definitions, sample, estimators, and test list before estimating. \\
4 & Analysis & Main event study and heterogeneity. \\
5 & Robustness & Placebos, alternative estimators and treatment definitions, sensitivity. \\
6 & Reporting & Final report, figures, replication notes. \\
\bottomrule
\end{tabular}

\subsection*{Proposed go/no-go criteria (after the data audit)}
\textbf{Go} if all hold: (i) enough treated users with usable pre- and post-periods that the minimum detectable effect is
economically meaningful; (ii) wearables and running products can be classified credibly; (iii) pre-trends in a first
look are not obviously violated. \textbf{Adjust} if only the outcome or treatment definition fails (redefine, re-audit once).
\textbf{No-go} if (i) fails, then switch data or reframe as descriptive. Numeric thresholds are left to the
decision owner, to be set \emph{before} looking at estimates.

\section{Open decisions}
\begin{tabular}{p{0.38\textwidth}p{0.30\textwidth}p{0.10\textwidth}p{0.10\textwidth}}
\toprule
Decision & Proposed default & Owner & Deadline \\
\midrule
Treatment definition & Narrow (trackers, GPS watches) & TBD & TBD \\
Primary outcome & Running-related purchases & TBD & TBD \\
Primary estimator & Callaway--Sant'Anna & TBD & TBD \\
Go/no-go thresholds & Set numerically before estimation & TBD & TBD \\
Final deliverable & Report series; paper optional & TBD & TBD \\
\bottomrule
\end{tabular}

\section{Next steps and follow-up issues}
\begin{enumerate}
  \item Review this proposal; resolve the open decisions above.
  \item Open one issue per phase: Setup, Data audit, Design, Analysis, Robustness, Reporting.
  \item Data audit issue includes the go/no-go memo as its deliverable.
\end{enumerate}

\end{document}
